\documentclass[11pt]{beamer}
\usepackage[utf8]{inputenc}
\usepackage[T1,T2A]{fontenc}
\usepackage[russian]{babel}
\usepackage{color}
\usepackage{listings}
\usepackage{calc}
\usepackage{graphicx}
\usepackage{epstopdf}
\usepackage{hyperref}
\hypersetup{unicode,colorlinks}
\usepackage{csquotes}
\usepackage{upquote}
\usepackage{cprotect}
\usetheme[progressbar=head,numbering=fraction,block=fill]{metropolis}
\usepackage{dejavu}
\usepackage{etoolbox}
\usepackage{bm}
\usepackage{minted}
%\usepackage[defaultsans]{droidsans}
%\usepackage[defaultmono]{droidmono}
%\makeatletter
%\input droid/t1fdm.fd
%\makeatother

\makeatletter
\DeclareRobustCommand*{\bfseries}{%
    \not@math@alphabet\bfseries\mathbf
    \fontseries\bfdefault\selectfont
    \boldmath
}
\makeatother

\makeatletter
\chardef\straightquote@code=\catcode`'
\chardef\backquote@code=\catcode``
\catcode`'=\active \catcode``=\active
\patchcmd{\@noligs}
{\textasciigrave}
{\fixedtextasciigrave}
{}{}
\newcommand{\fixedtextasciigrave}{%
    \makebox[.5em]{\fontencoding{TS1}\fontfamily{fvs}\selectfont\textasciigrave}% Vera Sans
}
\catcode\lq\'=\straightquote@code
\catcode\lq\`=\backquote@code
\makeatletter

\lstset{
    language=Haskell,
    inputencoding=utf8,
    extendedchars=true,
    breaklines=true,
    escapeinside=!!,
    tabsize=4,
    breakatwhitespace=true,
    keepspaces=true,
    basicstyle=\ttfamily,
    commentstyle=\color{gray}\itshape,
    keywordstyle=
    % keywordstyle=\color{blue}, % TODO https://tex.stackexchange.com/questions/415777/avoid-highlighting-keywords-following-certain-words-in-listings
    % addkeywords={Type}
    % deletekeywords={Monad, String, return, array, show},
    % otherkeywords={::, ->}
}
\lstset{
    literate={а}{{\selectfont\char224}}1
    {б}{{\selectfont\char225}}1
    {в}{{\selectfont\char226}}1
    {г}{{\selectfont\char227}}1
    {д}{{\selectfont\char228}}1
    {е}{{\selectfont\char229}}1
    {ё}{{\"e}}1
    {ж}{{\selectfont\char230}}1
    {з}{{\selectfont\char231}}1
    {и}{{\selectfont\char232}}1
    {й}{{\selectfont\char233}}1
    {к}{{\selectfont\char234}}1
    {л}{{\selectfont\char235}}1
    {м}{{\selectfont\char236}}1
    {н}{{\selectfont\char237}}1
    {о}{{\selectfont\char238}}1
    {п}{{\selectfont\char239}}1
    {р}{{\selectfont\char240}}1
    {с}{{\selectfont\char241}}1
    {т}{{\selectfont\char242}}1
    {у}{{\selectfont\char243}}1
    {ф}{{\selectfont\char244}}1
    {х}{{\selectfont\char245}}1
    {ц}{{\selectfont\char246}}1
    {ч}{{\selectfont\char247}}1
    {ш}{{\selectfont\char248}}1
    {щ}{{\selectfont\char249}}1
    {ъ}{{\selectfont\char250}}1
    {ы}{{\selectfont\char251}}1
    {ь}{{\selectfont\char252}}1
    {э}{{\selectfont\char253}}1
    {ю}{{\selectfont\char254}}1
    {я}{{\selectfont\char255}}1
    {А}{{\selectfont\char192}}1
    {Б}{{\selectfont\char193}}1
    {В}{{\selectfont\char194}}1
    {Г}{{\selectfont\char195}}1
    {Д}{{\selectfont\char196}}1
    {Е}{{\selectfont\char197}}1
    {Ё}{{\"E}}1
    {Ж}{{\selectfont\char198}}1
    {З}{{\selectfont\char199}}1
    {И}{{\selectfont\char200}}1
    {Й}{{\selectfont\char201}}1
    {К}{{\selectfont\char202}}1
    {Л}{{\selectfont\char203}}1
    {М}{{\selectfont\char204}}1
    {Н}{{\selectfont\char205}}1
    {О}{{\selectfont\char206}}1
    {П}{{\selectfont\char207}}1
    {Р}{{\selectfont\char208}}1
    {С}{{\selectfont\char209}}1
    {Т}{{\selectfont\char210}}1
    {У}{{\selectfont\char211}}1
    {Ф}{{\selectfont\char212}}1
    {Х}{{\selectfont\char213}}1
    {Ц}{{\selectfont\char214}}1
    {Ч}{{\selectfont\char215}}1
    {Ш}{{\selectfont\char216}}1
    {Щ}{{\selectfont\char217}}1
    {Ъ}{{\selectfont\char218}}1
    {Ы}{{\selectfont\char219}}1
    {Ь}{{\selectfont\char220}}1
    {Э}{{\selectfont\char221}}1
    {Ю}{{\selectfont\char222}}1
    {Я}{{\selectfont\char223}}1
    {`}{{{\fixedtextasciigrave}}}1
}

% \usemintedstyle{solarized-light}
\newminted[haskell]{haskell}{
    escapeinside=``,
    beameroverlays=true,
    texcomments=true,
    autogobble=true,
    fontsize=\small
}
\newmintinline[haskinline]{haskell}{
    escapeinside=``,
    beameroverlays=true,
    texcomments=true
}

\author{Алексей Романов}
\subtitle{Функциональное программирование на Haskell}
%\logo{}
\institute{МИЭТ}
\subject{Функциональное программирование на Haskell}
%\setbeamercovered{transparent}
%\setbeamertemplate{navigation symbols}{}

\usepackage{tikz}

\title{Лекция 6а: персистентные структуры данных}

\begin{document}
\begin{frame}[plain]
  \maketitle
\end{frame}

\begin{frame}[fragile]
  \frametitle{Две версии состояния}
  В SimpleLang из лабораторной 2 значение неприсвоенной переменной равно нулю.
  \begin{haskell}
    type State = String -> Integer

    s0 = empty
    s1 = extend s0 "x" 3
    s2 = extend s1 "x" 8
    s3 = extend s1 "y" 5
  \end{haskell}
  Чему равны \haskinline|s1 "x"|, \haskinline|s2 "x"| и \haskinline|s3 "x"|? \pause

  Ответ: 3, 8 и 3. От одной версии \haskinline|s1| можно получить две независимые версии. Нужно ли для этого дважды копировать всё состояние?
\end{frame}

\begin{frame}[fragile]
  \frametitle{Что такое персистентность?}
  \begin{itemize}
    \item Обновление возвращает новую версию структуры. Старой можно продолжать пользоваться, пока на неё есть ссылка.
    \item Если обновлять можно любую версию, история может ветвиться, как в примере с \haskinline|s1|.
    \item Как сделать это дёшево? \pause{} Старое и новое значения могут ссылаться на одни и те же неизменяемые части структуры. Это называют \emph{структурным разделением} (\emph{structural sharing}).
  \end{itemize}
\end{frame}

\begin{frame}[fragile]
  \frametitle{Знакомый пример: добавить голову}
  \begin{haskell}
    xs = [2, 3]
    ys = 1 : xs
    zs = 0 : xs
  \end{haskell}
  Сколько узлов списка добавила каждая операция \haskinline|(:)|? \pause

  \begin{center}
  \begin{tikzpicture}[x=1.15cm,y=.55cm,every node/.style={font=\small}]
    \node (y) at (0,1) {\texttt{ys}};
    \node (z) at (0,-1) {\texttt{zs}};
    \node[draw] (one) at (1,1) {1};
    \node[draw] (zero) at (1,-1) {0};
    \node (x) at (2.4,1.35) {\texttt{xs}};
    \node[draw] (two) at (2.4,0) {2};
    \node[draw] (three) at (3.5,0) {3};
    \node (nil) at (4.5,0) {\texttt{[]}};
    \foreach \a/\b in {y/one,z/zero,one/two,zero/two,x/two,two/three,three/nil}
      \draw[->] (\a) -- (\b);
  \end{tikzpicture}
  \end{center}
  По одному узлу. Хвост у всех трёх списков общий.
\end{frame}

\begin{frame}[fragile]
  \frametitle{А если изменить середину списка?}
  \begin{haskell}
    replaceFirst :: Eq a => a -> a -> [a] -> [a]
    replaceFirst _ _ [] = []
    replaceFirst old new (x : xs)
      | x == old  = new : xs
      | otherwise = x : replaceFirst old new xs
  \end{haskell}
  Для \haskinline|replaceFirst 3 9 [1,2,3,4,5]| какая часть результата может быть общей с исходным списком? \pause

  Хвост \haskinline|[4,5]|. Узлы с 1 и 2 в результате ведут к новому узлу с 9. При построении результата заново создаётся весь префикс до замены включительно.
\end{frame}

\begin{frame}[fragile]
  \frametitle{Множество на дереве поиска}
  \begin{haskell}
    data BST a = Empty | Branch (BST a) a (BST a)
  \end{haskell}
  В каждом узле все ключи слева меньше ключа узла, а справа "--- больше. Равные ключи храним один раз.
  \begin{haskellsmall}
    member :: Ord a => a -> BST a -> Bool
    member _ Empty = False
    member x (Branch l k r) =!\pause! case compare x k of
      LT -> member x l
      EQ -> True
      GT -> member x r
  \end{haskellsmall}
  Инвариант позволяет на каждом шаге выбрать одно поддерево.
\end{frame}

\begin{frame}[fragile]
  \frametitle{Вставка возвращает новое дерево}
  Как при возврате из рекурсии собрать путь от корня до нового ключа?
  \begin{haskell}
    insert :: Ord a => a -> BST a -> BST a
    insert x Empty = Branch Empty x Empty
    insert x t@(Branch l k r) =!\pause! case compare x k of
      LT -> Branch (insert x l) k r
      EQ -> t
      GT -> Branch l k (insert x r)
  \end{haskell}
  В ветке \haskinline|LT| поле \haskinline|r| переходит в новый узел целиком. Что переходит целиком в ветке \haskinline|GT|?
\end{frame}

\begin{frame}[fragile]
  \frametitle{Куда указывает новый корень?}
  Вставим 5 в \haskinline|t|. Какие узлы сохранит \haskinline|insert 5 t|?
  \begin{center}
  \begin{tikzpicture}[x=.85cm,y=.76cm,vertex/.style={circle,draw,inner sep=2pt,minimum size=5.5mm},new/.style={vertex,very thick},every node/.style={font=\small}]
    \node (t) at (0,1.7) {\texttt{t}};
    \node[vertex] (a) at (0,.9) {8};
    \node[vertex] (b) at (-1.6,0) {3};
    \node[vertex] (c) at (5.6,0) {10};
    \node[vertex] (d) at (-1.6,-1.1) {1};
    \foreach \i/\j in {t/a,a/b,b/d} \draw[->] (\i)--(\j);
    \draw[->] (a) to[out=18,in=80] (c);
    \uncover<2->{
      \node (u) at (4,1.7) {\texttt{insert 5 t}};
      \node[new] (e) at (4,.9) {8};
      \node[new] (f) at (2.6,.1) {3};
      \node[new] (g) at (3.7,-1.1) {5};
      \foreach \i/\j in {u/e,e/f,f/g} \draw[->] (\i)--(\j);
      \draw[->] (e) -- (c);
      \draw[->] (f) to[bend left=30] (d);
    }
  \end{tikzpicture}
  \end{center}
  \uncover<2->{В новой версии три новых узла: 8, 3 и 5. Обе версии ссылаются на прежние узлы 1 и 10. Пустые поддеревья на рисунке опущены.}
\end{frame}

\begin{frame}[fragile]
  \frametitle{Сколько стоит обновление?}
  \begin{itemize}
    \item Пусть $h$ "--- высота дерева. При вставке нового ключа посещаем один путь: $O(h+1)$ сравнений и столько же новых узлов в худшем случае.
    \item Какую форму примет дерево при вставке ключей 1, 2, 3, \ldots{} по порядку? \pause
    \item Каждый следующий ключ уйдёт вправо. Высота может вырасти до $O(n)$ для $n$ ключей.
    \item Сбалансированное дерево поддерживает высоту $O(\log n)$. В библиотечных множествах и словарях балансировка входит в реализацию.
  \end{itemize}
\end{frame}

\begin{frame}[fragile]
  \frametitle{Что остаётся в памяти?}
  \begin{itemize}
    \item Если сохранить только новый корень, что произойдёт со старым путём от корня до места вставки? \pause
    \item Сборщик мусора может освободить узлы, до которых уже нельзя добраться. Общие поддеревья остаются доступны из нового корня.
    \item Если сохранять все версии, их корни удерживают и старые пути. История из $m$ вставок в сбалансированное дерево размера не более $n$ требует до $O(m\log(n+1))$ новых узлов сверх исходного дерева.
  \end{itemize}
\end{frame}

\begin{frame}[fragile]
  \frametitle{Множество в библиотеке}
  Ключи без повторений хранятся в \href{https://hackage.haskell.org/package/containers/docs/Data-Set.html}{\haskinline{Data.Set}} из пакета \hackage{containers}.
  \begin{haskell}
    import qualified Data.Set as S

    used = S.fromList ["x", "y", "x"]
    more = S.insert "z" used
  \end{haskell}
  Что вернёт выражение?
  \begin{haskellsmall}
    ghci> (S.size used, S.member "z" used, S.size more)
  \end{haskellsmall}
  \pause
  \begin{haskellsmall}
    (2,False,3)
  \end{haskellsmall}
  Для поиска и вставки достаточно $O(\log n)$ сравнений.
\end{frame}

\begin{frame}[fragile]
  \frametitle{Словарь: ключ и значение}
  \href{https://hackage.haskell.org/package/containers/docs/Data-Map-Strict.html}{\haskinline{Data.Map.Strict}} сопоставляет каждому ключу одно значение.
  \begin{haskellsmall}
    import qualified Data.Map.Strict as M

    env1 = M.singleton "x" 3
    env2 = M.insert "x" 8 env1
  \end{haskellsmall}
  \begin{haskellsmall}
    ghci> (M.lookup "x" env1, M.lookup "x" env2)
  \end{haskellsmall}
  \pause
  \begin{haskellsmall}
    (Just 3,Just 8)
    ghci> M.lookup "y" env2
  \end{haskellsmall}
  \pause
  \begin{haskellsmall}
    Nothing
  \end{haskellsmall}
  \haskinline|M.lookup :: Ord k => k -> M.Map k v -> Maybe v|
\end{frame}

\begin{frame}[fragile]
  \frametitle{Пример: частоты слов}
  Построим словарь, в котором ключ "--- слово, а значение "--- число его появлений.
  \begin{haskellsmall}
    counts :: [String] -> M.Map String Int
    counts = foldl' add M.empty
      where
        add m word = M.insertWith (+) word 1 m
  \end{haskellsmall}
  \haskinline|foldl'| "--- строгий вариант \haskinline|foldl|; подробнее в лекции про ленивость.

  \haskinline|insertWith| при совпадении ключа вычисляет \haskinline|combine new old|. Для нового ключа записывает переданное значение.
  \begin{haskellsmall}
    ghci> M.toAscList (counts ["cat", "dog", "cat"])
  \end{haskellsmall}
  \pause
  \begin{haskellsmall}
    [("cat",2),("dog",1)]
  \end{haskellsmall}
\end{frame}

\begin{frame}[fragile]
  \frametitle{Что выбрать?}
  \begin{itemize}
    \item Нужны все элементы в порядке поступления, включая повторы: список. Если добавляем в начало, \haskinline|reverse| в конце восстановит порядок поступления.
    \item Нужно быстро проверять принадлежность и хранить каждый ключ один раз: \haskinline|Set a|.
    \item Нужно получать значение по ключу: \haskinline|Map k v|.
  \end{itemize}
  Упражнение: в каком типе хранить все вхождения слов по порядку? Уже встреченные слова? Число вхождений каждого слова? \pause

  \haskinline|[String]|, \haskinline|S.Set String| и \haskinline|M.Map String Int| соответственно.
\end{frame}

\begin{frame}[fragile]
  \frametitle{Вернёмся к SimpleLang}
  Реализуйте состояние на основе словаря:
  \begin{haskellsmall}
    type Env = M.Map String Integer

    emptyEnv :: Env
    readVar  :: Env -> String -> Integer
    assign   :: Env -> String -> Integer -> Env
  \end{haskellsmall}
  \begin{itemize}
    \item Как читать неприсвоенную переменную? По правилам SimpleLang её значение равно нулю.
    \item Повторите пример со слайда 2. Проверьте результат чтения в каждой ветви.
  \end{itemize}
\end{frame}

\begin{frame}[fragile]
  \frametitle{Для самостоятельной работы}
  \begin{enumerate}
    \item Нарисуйте ещё одну версию дерева со слайда про вставку: добавьте 12 в \emph{исходное} дерево. Какие узлы общие у всех трёх версий?
    \item Запишите свойство, связывающее \haskinline|member| и \haskinline|insert|. Оно должно описывать поиск как вставленного ключа, так и любого другого.
    \item Отладчик SimpleLang хранит список предыдущих состояний. Как добавить шаг и как отменить последний шаг? Если оставить только последнее состояние, что останется в памяти? Сравните \haskinline|Env| на основе словаря и \haskinline|State| из лабораторной 2, где новая функция обращается к предыдущей.
  \end{enumerate}
\end{frame}

\begin{frame}
  \frametitle{Дополнительное чтение}
  \begin{itemize}
    \item \href{https://haskell-containers.readthedocs.io/en/latest/}{Руководство по containers}. Разделы \emph{Sets} и \emph{Maps}: выбор операций, примеры и оценки сложности.
    \item Chris Okasaki, \href{https://www.cs.cmu.edu/~rwh/students/okasaki.pdf}{Purely Functional Data Structures}, глава 2 диссертации. Списки и деревья поиска как примеры сохранения старых версий.
  \end{itemize}
\end{frame}
\end{document}